\begin{abstract}
We introduce \ours, a remarkably simple modification to standard deep
ensembling: train the same architecture with the same random seed and the
same cosine learning-rate schedule for two different lengths (e.g., 100 vs.\
150 epochs), and ensemble the two resulting checkpoints. The two training
trajectories reach \emph{identical} mean validation accuracy (differing by
$<\!0.001$ in our experiments), but they differ on 13\% of test samples
because the longer trajectory spends 50 additional epochs in a low-LR
``basin-exploration'' phase that drifts weights into a neighboring local
minimum. Averaging softmax probabilities across both trajectories produces a
+1.5 percentage-point gain \emph{beyond} what standard seed-only deep
ensembling achieves.

We validate \ours across two modalities and six datasets. In EEG emotion
classification, we achieve a new state of the art on FACED 9-class
(\facedacc{} balanced accuracy, +5.3 pp / +8.2\% over the EMOD AAAI 2026
baseline), matched by a cross-dataset confirmation on SEED-V (BAcc 0.4083,
\seedvgain{} ensemble gain). In image classification, we replicate the
trick on CIFAR-10 (accuracy \ciftenacc{}, \ciftengain{} gain), CIFAR-100,
MNIST, Fashion-MNIST, and SVHN. Head-to-head against SWA, snapshot
ensembles, and standard 10-seed deep ensembles, \ours{} is competitive or
stronger while adding \emph{zero} hyperparameter tuning and a single
additional training run.

A theoretical decomposition shows that \ours-induced diversity is
proportional to the ``unused'' individual-accuracy headroom: gain scales from
\facedgain{} on FACED (individual mean 0.66) down to +0.14\% on MNIST
(individual mean 0.985). The trick is most valuable where ensembling is
most needed---on capacity-bound students near their effective ceiling.
\end{abstract}
